\chapter{Implementation}
\label{ch:implementation}

This chapter describes the practical implementation of the Lambda cold-start isolation experiment, including packaging procedures, handler implementations, log parsing mechanisms, and statistical test implementations.

\section{Lambda Function Handlers}

Each runtime implements an identical handler interface: receive JSON payload, execute SHA-256 computation, return structured response.

\subsection{Python Handler}

\textbf{File:} \texttt{functions/python/default/handler.py} (and \texttt{.../optimised/handler.py}, identical code)

\begin{verbatim}
import json
import hashlib

# Default variant imports (never used):
# import boto3
# import requests
# import sympy

def lambda_handler(event, context):
    if event.get("warmer"):
        return {"ok": True, "warmer": True}
    
    n = event["n"]
    items = event["items"]
    
    digest = "chaitanya-25171216-nci"
    for _ in range(20000):
        digest = hashlib.sha256(digest.encode()).hexdigest()
    
    return {
        "ok": True,
        "digest": digest,
        "n": n,
        "items": items,
        "items_total": sum(items)
    }
\end{verbatim}

The \texttt{default} variant includes import statements for boto3, requests, and sympy at the top of the file, but these libraries are never referenced in the code. This simulates a real-world scenario where dependencies accumulate over time but are not removed.

Packaging (\texttt{scripts/package\_all.sh}):
\begin{verbatim}
pip install boto3 requests sympy -t functions/python/default/
zip -r python-default.zip functions/python/default/
\end{verbatim}

The \texttt{optimised} variant has no \texttt{pip install} step, resulting in a 604-byte ZIP containing only the handler script.

\subsection{Node.js Handler}

\textbf{File:} \texttt{functions/nodejs/default/index.mjs}

\begin{verbatim}
import crypto from 'crypto';

// Default variant requires (never called):
// import AWS from 'aws-sdk';
// import _ from 'lodash';
// import moment from 'moment';
// import axios from 'axios';

export const handler = async (event) => {
    if (event.warmer) {
        return {ok: true, warmer: true};
    }
    
    const {n, items} = event;
    
    let digest = "chaitanya-25171216-nci";
    for (let i = 0; i < 20000; i++) {
        digest = crypto.createHash('sha256')
                       .update(digest)
                       .digest('hex');
    }
    
    return {
        ok: true,
        digest,
        n,
        items,
        items_total: items.reduce((a, b) => a + b, 0)
    };
};
\end{verbatim}

The \texttt{default} variant includes a \texttt{package.json} with dependencies:
\begin{verbatim}
{
  "type": "module",
  "dependencies": {
    "aws-sdk": "^2.1691.0",
    "lodash": "^4.17.21",
    "moment": "^2.30.1",
    "axios": "^1.7.9"
  }
}
\end{verbatim}

Packaging: \texttt{npm install --production} (16.5 MB ZIP) vs. no npm install for optimised (815-byte ZIP).

\subsection{Java Handler}

\textbf{File:} \texttt{functions/java/default/src/main/java/coldstart/Handler.java}

\begin{verbatim}
package coldstart;

import com.amazonaws.services.lambda.runtime.Context;
import com.amazonaws.services.lambda.runtime.RequestHandler;
import java.security.MessageDigest;
import java.util.Map;
import java.util.List;

// Default variant touches (never meaningfully used):
// import com.fasterxml.jackson.databind.ObjectMapper;
// import com.google.common.base.Strings;
// import org.apache.commons.lang3.StringUtils;
// static { ObjectMapper om = new ObjectMapper(); }

public class Handler implements RequestHandler<Map, Map> {
    public Map handleRequest(Map event, Context context) {
        if (event.containsKey("warmer")) {
            return Map.of("ok", true, "warmer", true);
        }
        
        int n = (int) event.get("n");
        List<Integer> items = (List<Integer>) event.get("items");
        
        String digest = "chaitanya-25171216-nci";
        try {
            MessageDigest md = MessageDigest.getInstance("SHA-256");
            for (int i = 0; i < 20000; i++) {
                digest = bytesToHex(md.digest(digest.getBytes()));
            }
        } catch (Exception e) {
            throw new RuntimeException(e);
        }
        
        int sum = items.stream().mapToInt(Integer::intValue).sum();
        
        return Map.of(
            "ok", true,
            "digest", digest,
            "n", n,
            "items", items,
            "items_total", sum
        );
    }
    
    private static String bytesToHex(byte[] bytes) {
        StringBuilder sb = new StringBuilder();
        for (byte b : bytes) {
            sb.append(String.format("%02x", b));
        }
        return sb.toString();
    }
}
\end{verbatim}

The \texttt{default} variant includes a \texttt{pom.xml} with shaded dependencies:
\begin{verbatim}
<dependencies>
  <dependency>
    <groupId>com.fasterxml.jackson.core</groupId>
    <artifactId>jackson-databind</artifactId>
    <version>2.18.2</version>
  </dependency>
  <dependency>
    <groupId>com.google.guava</groupId>
    <artifactId>guava</artifactId>
    <version>33.4.0-jre</version>
  </dependency>
  <dependency>
    <groupId>org.apache.commons</groupId>
    <artifactId>commons-lang3</artifactId>
    <version>3.17.0</version>
  </dependency>
</dependencies>
\end{verbatim}

Maven build with \texttt{maven-shade-plugin} produces a 6.3 MB JAR. The \texttt{optimised} variant has no dependencies and produces a 4.9 KB JAR.

\section{Log Parsing Implementation}

\subsection{REPORT Line Format}

AWS Lambda logs execution metrics in a structured format. Example cold-start REPORT:

\begin{verbatim}
REPORT RequestId: 12345-abcde-67890
Duration: 823.45 ms
Billed Duration: 900 ms
Memory Size: 512 MB
Max Memory Used: 187 MB
Init Duration: 756.12 ms
\end{verbatim}

Warm (container reused):

\begin{verbatim}
REPORT RequestId: 67890-fghij-12345
Duration: 9.23 ms
Billed Duration: 100 ms
Memory Size: 512 MB
Max Memory Used: 188 MB
\end{verbatim}

\subsection{Parser Implementation}

\textbf{File:} \texttt{src/coldstart/report\_parser.py}

The parser uses compiled regex patterns for efficiency:

\begin{verbatim}
import re

REPORT_PATTERN = re.compile(
    r'REPORT RequestId: (?P<request_id>[\w-]+).*?'
    r'Duration: (?P<duration>[\d.]+) ms.*?'
    r'Billed Duration: (?P<billed>[\d]+) ms.*?'
    r'Memory Size: (?P<memory>[\d]+) MB.*?'
    r'Max Memory Used: (?P<max_memory>[\d]+) MB'
    r'(?:.*?Init Duration: (?P<init>[\d.]+) ms)?',
    re.DOTALL
)

def parse_report_line(line):
    match = REPORT_PATTERN.search(line)
    if not match:
        return None
    
    d = match.groupdict()
    init_duration = float(d["init"]) if d["init"] else None
    
    return {
        "request_id": d["request_id"],
        "duration_ms": float(d["duration"]),
        "billed_duration_ms": int(d["billed"]),
        "memory_size_mb": int(d["memory"]),
        "max_memory_used_mb": int(d["max_memory"]),
        "init_duration_ms": init_duration,
        "is_cold": init_duration is not None
    }
\end{verbatim}

The \texttt{(?:...)?} non-capturing optional group matches Init Duration only when present. This aligns with AWS's official cold-start definition.

\section{Statistical Test Implementation}

\subsection{Normality Testing}

\textbf{File:} \texttt{src/coldstart/stats.py}

\begin{verbatim}
from scipy.stats import shapiro

def test_normality(samples, alpha=0.05):
    """
    Shapiro-Wilk test.
    Returns: (is_normal: bool, p_value: float)
    """
    if len(samples) < 3:
        return False, None  # Too few samples
    
    stat, p = shapiro(samples)
    return p > alpha, p
\end{verbatim}

\subsection{Hypothesis Test Selection}

Based on normality results, the analysis selects parametric or non-parametric tests:

\begin{verbatim}
from scipy.stats import (
    kruskal, mannwhitneyu, wilcoxon,
    f_oneway, ttest_ind, ttest_rel
)

def compare_groups(groups, paired=False):
    """
    Groups: dict of {label: [samples]}
    Returns: test name, statistic, p-value, effect size
    """
    # Check normality for each group
    all_normal = all(
        test_normality(samples)[0] 
        for samples in groups.values()
    )
    
    if len(groups) == 2:
        labels = list(groups.keys())
        a, b = groups[labels[0]], groups[labels[1]]
        
        if paired:
            if all_normal:
                stat, p = ttest_rel(a, b)
                effect = cohens_d_paired(a, b)
                return "paired_t", stat, p, effect
            else:
                stat, p = wilcoxon(a, b)
                effect = rank_biserial(a, b)
                return "wilcoxon", stat, p, effect
        else:
            if all_normal:
                stat, p = ttest_ind(a, b, equal_var=False)
                effect = cohens_d(a, b)
                return "welch_t", stat, p, effect
            else:
                stat, p = mannwhitneyu(a, b)
                effect = rank_biserial(a, b)
                return "mann_whitney_u", stat, p, effect
    
    else:  # Multiple groups
        if all_normal:
            stat, p = f_oneway(*groups.values())
            effect = eta_squared(groups)
            return "anova", stat, p, effect
        else:
            stat, p = kruskal(*groups.values())
            effect = epsilon_squared(groups, stat)
            return "kruskal_wallis", stat, p, effect
\end{verbatim}

\subsection{Multiple Comparison Correction}

Holm-Bonferroni sequential correction:

\begin{verbatim}
def holm_bonferroni(p_values, alpha=0.05):
    """
    p_values: list of (label, p_value) tuples
    Returns: list of (label, p_original, p_corrected, reject)
    """
    sorted_p = sorted(p_values, key=lambda x: x[1])
    m = len(sorted_p)
    
    results = []
    for i, (label, p) in enumerate(sorted_p):
        corrected_alpha = alpha / (m - i)
        reject = p <= corrected_alpha
        results.append((label, p, corrected_alpha, reject))
    
    return results
\end{verbatim}

\section{Cost Model Implementation}

\textbf{File:} \texttt{src/coldstart/cost\_model.py}

AWS Lambda pricing (US East, September 2026):

\begin{itemize}
    \item Request charge: \$0.20 per 1M requests = \$0.0000002 per request
    \item Duration charge: \$0.0000166667 per GB-second
\end{itemize}

\begin{verbatim}
def compute_invocation_cost(memory_mb, billed_duration_ms):
    """
    Returns: cost in USD
    """
    request_cost = 0.0000002
    
    memory_gb = memory_mb / 1024.0
    duration_sec = billed_duration_ms / 1000.0
    gb_sec = memory_gb * duration_sec
    duration_cost = gb_sec * 0.0000166667
    
    return request_cost + duration_cost

def cost_per_1k(memory_mb, billed_duration_ms):
    """
    Returns: cost per 1,000 invocations in USD
    """
    return 1000 * compute_invocation_cost(
        memory_mb, billed_duration_ms
    )
\end{verbatim}

\subsection{Decision Matrix Construction}

The decision matrix expresses each optimization as milliseconds saved per additional dollar per 1,000 invocations:

\begin{verbatim}
def decision_matrix_row(
    baseline_latency_ms, 
    baseline_cost_1k_usd,
    optimized_latency_ms, 
    optimized_cost_1k_usd
):
    """
    Returns: {
        "latency_reduction_ms": float,
        "cost_increase_usd": float,
        "roi_ms_per_dollar": float,
        "recommendation": str
    }
    """
    delta_latency = baseline_latency_ms - optimized_latency_ms
    delta_cost = optimized_cost_1k_usd - baseline_cost_1k_usd
    
    if delta_cost == 0:
        roi = float('inf') if delta_latency > 0 else 0
    else:
        roi = delta_latency / delta_cost
    
    # Recommendation bands
    if roi > 100:
        rec = "ADOPT"
    elif roi > 10:
        rec = "Consider"
    elif roi > 0:
        rec = "Marginal"
    else:
        rec = "Avoid"
    
    return {
        "latency_reduction_ms": delta_latency,
        "cost_increase_usd": delta_cost,
        "roi_ms_per_dollar": roi,
        "recommendation": rec
    }
\end{verbatim}

\section{Mock Backend Implementation}

The mock backend generates synthetic data for pipeline validation.

\textbf{File:} \texttt{src/coldstart/mock.py}

\begin{verbatim}
import numpy as np

class MockBackend:
    def __init__(self, config):
        self.runtimes = config["runtimes"]
        self.variance_cv = config["variance_cv"]
    
    def generate_init_duration(
        self, 
        runtime, 
        package_variant, 
        memory_mb
    ):
        """
        Returns: init_duration_ms (float)
        """
        base = self.runtimes[runtime]["init_base_ms"]
        
        if package_variant == "default":
            base += self.runtimes[runtime]["package_penalty_ms"]
        
        # Memory effect (exploratory H4)
        if memory_mb >= 1769:
            base *= 0.85  # vCPU coupling speedup
        
        # Add variance
        std = base * self.variance_cv
        sample = np.random.normal(base, std)
        
        return max(sample, 10.0)  # Clamp to realistic minimum
\end{verbatim}

Configuration (\texttt{configs/mock\_model.yaml}):

\begin{verbatim}
runtimes:
  python:
    init_base_ms: 800
    package_penalty_ms: 2000
  nodejs:
    init_base_ms: 300
    package_penalty_ms: 280
  java:
    init_base_ms: 250
    package_penalty_ms: 180

variance_cv: 0.15
\end{verbatim}

\section{Deployment Automation}

\textbf{File:} \texttt{scripts/deploy.sh}

\begin{verbatim}
#!/bin/bash
set -e

echo "Building packages..."
bash scripts/package_all.sh python nodejs java

echo "Validating digest uniformity..."
python3 scripts/check_digests.py

echo "SAM build..."
cd infra
sam build

echo "SAM deploy..."
sam deploy --no-confirm-changeset

echo "Deployment complete. Function ARNs:"
aws cloudformation describe-stacks \
    --stack-name lambda-coldstart-isolation \
    --query 'Stacks[0].Outputs' \
    --output table
\end{verbatim}

\section{Summary}

This chapter detailed the practical implementation of Lambda handlers, log parsing, statistical tests, cost modeling, and deployment automation. Key implementation decisions include:

\begin{itemize}
    \item Identical handler logic across runtimes (only packaging differs)
    \item Regex-based REPORT line parsing with explicit cold-start detection (Init Duration presence)
    \item Automated test selection based on Shapiro-Wilk normality
    \item AWS pricing-accurate cost model with decision matrix ROI calculation
    \item Mock backend for zero-cost pipeline validation
\end{itemize}

The next chapter presents experimental results, statistical analysis outcomes, and decision matrix construction (contingent on live AWS execution).
